% LSZ 条件、约化公式和独立证明在正文中相邻排列。
实验真正制备和探测的是进入相互作用区之前、离开相互作用区之后能够独立传播的粒子；而场论计算首先给出的却是若干时空点上的场关联函数。两者之间因此还缺少一步：必须从相关函数中取出与一个真实外部粒子对应的传播部分，只留下相互作用本身的振幅。

对稳定粒子，这一步可以从精确二点函数的质量壳结构看清楚。当四动量接近真实粒子的质量壳时，二点函数出现一个单粒子极点；极点的位置给出物理质量，极点的留数给出场算符与规范化单粒子态之间的重叠。LSZ 约化所做的就是用自由逆运动算符消去这段外部传播，再按极点留数归一化，并把结果投影到真实电子自旋态或真实光子的横向偏振态。得到的剩余部分才是散射振幅。

因此，“截肢外腿”不是在图上把最外面的一条线擦掉，而是对质量壳极点执行一个明确的代数留数运算。内部传播线并不代表渐近探测到的单粒子态，所以不进行同样的质量壳投影。


光子外腿使用物理四动量 $k^\mu=\hbar\kappa^\mu$。把\eqnrefs{eq:maxwell-kernel-projectors-r6} 的四波矢核改写成物理四动量核，定义
\begin{equation}
 \mathscr K_{\xi_{\rm gf}}^{\mu\nu}(k)
 \equiv
 \mathcal K_{\xi_{\rm gf}}^{\mu\nu}(\kappa=k/\hbar)
 =-\frac{k^2}{\mu_0\hbar^2}
 \left(P_T^{\mu\nu}+\frac1{\xi_{\rm gf}}P_L^{\mu\nu}\right).
 \label{eq:photon-physical-momentum-inverse-kernel-dp42}
\end{equation}
由\eqnrefs{eq:photon-prop-kappa-to-q-r13}，自由核与物理四动量表示的二点函数满足 $\mathscr K_{\mu\rho}\widetilde D_F^{\rho\nu}=\ii\hbar\,\delta_\mu{}^\nu$。电子传播子同样满足自由 Dirac 逆核乘传播子等于 $\ii\hbar$。因此两类外腿的截肢算符都必须含 $(\ii\hbar)^{-1}$：
\begin{equation}
 \boxed{
 \begin{aligned}
 \text{电子外腿:}&\quad
 \widetilde S(p)\ \xrightarrow{\;[Z_{\rm pole}^{(e)}]^{-1/2}(\slashed p-mc)/(\ii\hbar)\;}\
 u,\bar u,v,\bar v,\\
 \text{光子外腿:}&\quad
 \widetilde D_{\mu\nu}(k)\
 \xrightarrow{\;[Z_{\rm pole}^{(\gamma)}]^{-1/2}\mathscr K^{\mu\rho}(k)/(\ii\hbar)\;}\
 \varepsilon_\rho^{(\lambda)}(k),
 \qquad k\cdot\varepsilon^{(\lambda)}=0.
 \end{aligned}}
 \label{eq:lsz-dirac-photon-mainline-dp8}
\end{equation}
这里用 $Z_{\rm pole}^{(e)}$ 与 $Z_{\rm pole}^{(\gamma)}$ 专门表示当前所用场变量的电子与横向光子单粒子极点权重。它们是二点函数的谱数据，负责把局域场算符产生的状态归一化成单位范数渐近粒子。重整化中常用的 $Z_2,Z_3$ 则表示“裸场与重整化场之间的场强重整化常数”；在一般方案中，它们与单粒子极点权重是不同对象。若采用在壳方案，$Z_2,Z_3$ 被选择来使重整化传播子的相应物理极点留数等于 $1$。

\begin{figure}[htbp]
 \centering
 \begin{tikzfig}[x=1cm,y=1cm,>=Stealth]
 \node[font=\bfseries] at (-3.5,1.55) {(a) full Green function};
 \node[font=\bfseries] at (3.55,1.55) {(b) amputated pole};

 % electron row
 \draw[fermion] (-5.3,.55)--(-2.7,.55);
 \node[above] at (-4.0,.55) {$S(p)$};
 \node[draw,rounded corners=3pt,minimum width=1.45cm,minimum height=.9cm] at (-1.95,.55) {$\Gamma$};
 \node at (0,.55) {$\xrightarrow{\quad Z_2^{-1/2}(\slashed p-mc)\quad}$};
 \node at (2.15,.55) {$\bar u_s(p)$};
 \draw[fermion] (2.75,.55)--(3.35,.55);
 \node[draw,rounded corners=3pt,minimum width=1.45cm,minimum height=.9cm] at (4.1,.55) {$\Gamma$};

 % photon row
 \draw[photon] (-5.3,-1.0)--(-2.7,-1.0);
 \node[above] at (-4.0,-1.0) {$D_{\mu\nu}(k)$};
 \node[draw,rounded corners=3pt,minimum width=1.45cm,minimum height=.9cm] at (-1.95,-1.0) {$\Gamma^\nu$};
 \node at (0,-1.0) {$\xrightarrow{\quad Z_3^{-1/2}\mathcal K^{\rho\mu}_\xi\quad}$};
 \node at (2.05,-1.0) {$\varepsilon_\rho^{(\lambda)*}$};
 \draw[photon] (2.85,-1.0)--(3.35,-1.0);
 \node[draw,rounded corners=3pt,minimum width=1.45cm,minimum height=.9cm] at (4.1,-1.0) {$\Gamma^\rho$};
\end{tikzfig}

 \caption{LSZ 约化的图形表示。完整 Green 函数含外腿单粒子极点；自由逆运动核除以 $\ii\hbar$ 后消去该传播极点，再投影到壳上电子自旋子或横向光子偏振。$\Gamma$ 表示其余连通截肢核。}
 \label{fig:qed-lsz-amputation-dp8}
\end{figure}

\begin{proposition}[稳定 Dirac 单粒子极点的 LSZ 约化]
设精确二点函数有稳定、孤立的单粒子极点，$Z_{\rm pole}^{(e)}>0$。
连通 Green 函数满足 $\widetilde G=\widetilde S\widetilde\Gamma+\widetilde G_{\rm reg}$，原始截肢核与余项在极点邻域解析。
以下 $m=m_{\rm phys}$；限式先提取极点，再与对应质量壳动量的旋量收缩，
$u_s(p)$ 的记号表示壳上投影，并不假设离壳 $p$ 也有自由 Dirac 零模。严格无质量 QED 的带电红外粒子不自动满足孤立极点假设，必须另指定调节或适用的渐近态方案。
电子场与物理单电子态的重叠定义电子极点留数：
\begin{equation}
 \boxed{
 \langle\Omega|\psi(0)|p,s\rangle
 =\sqrt{Z_{\rm pole}^{(e)}}\,u_s(p),}
 \qquad
 (\slashed p-m_{\rm phys}c)u_s(p)=0.
 \label{eq:dirac-field-pole-overlap-dp8}
\end{equation}
于是精确电子二点函数在质量壳附近具有
\begin{equation}
 \boxed{
 \widetilde S(p)
 \xrightarrow[p^2\to m_{\rm phys}^2c^2]{}
 \frac{\ii\hbar Z_{\rm pole}^{(e)}(\slashed p+m_{\rm phys}c)}
 {p^2-m_{\rm phys}^2c^2+\ii0^+}
 +\widetilde S_{\rm reg}(p).}
 \label{eq:dirac-exact-pole-lsz-dp8}
\end{equation}
$\widetilde S_{\rm reg}$ 在该单粒子极点邻域解析，并包含多粒子连续谱贡献。若连通 Green 函数的一条出射电子外腿在该极点附近写成 $\widetilde G(p;\{q\})=\widetilde S(p)\widetilde\Gamma(p;\{q\})+\text{非极点项}$，则电子外腿的留数操作为
\begin{equation}
 \boxed{
 \lim_{p^2\to m_{\rm phys}^2c^2}
 \bar u_s(p)[Z_{\rm pole}^{(e)}]^{-1/2}
 \frac{\slashed p-m_{\rm phys}c}{\ii\hbar}
 \widetilde G(p;\{q\})
 =\sqrt{Z_{\rm pole}^{(e)}}\,\bar u_s(p)\widetilde\Gamma(p;\{q\}).}
 \label{eq:dirac-lsz-amputation-r68}
\end{equation}
\end{proposition}
\begin{proof}
设任意连通 Green 函数含一条出射 Dirac 外腿，写成
\begin{equation*}
 \widetilde G_\alpha(p;\{q\})
 =\int\dd^4 x\,\ee^{+\ii p\cdot x/\hbar}
 \langle\Omega|T\{\psi_\alpha(x)\mathcal O(\{q\})\}|\Omega\rangle_c.
\end{equation*}
若完整谱在 $p^2=m_{\rm phys}^2c^2$ 处含稳定单粒子极点，则极点邻域必可因子化为
\begin{equation*}
 \widetilde G(p;\{q\})
 =\widetilde S(p)\,\widetilde\Gamma(p;\{q\})
 +\widetilde G_{\rm reg}(p;\{q\}),
\end{equation*}
其中 $\widetilde\Gamma$ 与 $\widetilde G_{\rm reg}$ 在该极点解析。正文\eqnrefs{eq:dirac-field-pole-overlap-dp8} 给出场--粒子重叠
\begin{equation*}
 \langle\Omega|\psi(0)|p,s\rangle=\sqrt{Z_{\rm pole}^{(e)}}\,u_s(p),
\end{equation*}
所以精确二点函数的单粒子极点部分为
\begin{equation*}
 \widetilde S(p)
 \xrightarrow[p^2\to m_{\rm phys}^2c^2]{}
 \frac{\ii\hbar Z_{\rm pole}^{(e)}(\slashed p+m_{\rm phys}c)}
 {p^2-m_{\rm phys}^2c^2+\ii0^+}
 +\widetilde S_{\rm reg}(p).
\end{equation*}
左乘自由逆 Dirac 核并使用
\begin{equation*}
 (\slashed p-mc)(\slashed p+mc)=p^2-m^2c^2
\end{equation*}
可知
\begin{equation*}
 [Z_{\rm pole}^{(e)}]^{-1/2}\frac{\slashed p-mc}{\ii\hbar}\,\widetilde S(p)
 \xrightarrow[p^2\to m^2c^2]{}\sqrt{Z_{\rm pole}^{(e)}}
\end{equation*}
在取单粒子极点留数的意义下成立。于是
\begin{equation*}
 \lim_{p^2\to m^2c^2}
 \bar u_s(p)\,
 [Z_{\rm pole}^{(e)}]^{-1/2}\frac{\slashed p-mc}{\ii\hbar}
 \widetilde G(p;\{q\})
 =\sqrt{Z_{\rm pole}^{(e)}}\,\bar u_s(p)\widetilde\Gamma(p;\{q\}),
\end{equation*}
即正文\eqnrefs{eq:dirac-lsz-amputation-r68}。连续谱和其它在该点解析的部分没有相同极点，因此在取留数时消失；严格无质量 QED 的红外粒子结构需按正文给出的调节/包容方案另行处理。
原始截肢核因此还要乘 $\sqrt{Z_{\rm pole}^{(e)}}$，才对应规范化渐近态；其来源是 $Z\cdot Z^{-1/2}=\sqrt Z$。
这里的未投影逆核乘积只是在提取单粒子极点部分的意义下取极限，不能忽略完整矩阵的解析项。
最终与壳上 $\bar u$ 收缩时，逆 Dirac 核的壳上零模消去解析的非极点部分。

\end{proof}


这条式子把“二点函数具有单粒子极点”转换成“外腿可以被截肢并投影到壳上自旋子”。电子与正电子的四种渐近外态只是同一极点投影在交叉变换下的不同取向：
\begin{equation}
 \begin{array}{c|c}
 \text{渐近外态} & \text{壳上波函数}\\
 \hline
 \text{入射 }e^- & u_s(p)\\
 \text{出射 }e^- & \bar u_s(p)\\
 \text{入射 }e^+ & \bar v_s(p)\\
 \text{出射 }e^+ & v_s(p)
 \end{array}
 \label{eq:external-dirac-spinor-table-dp8}
\end{equation}

\begin{proposition}[横向光子单粒子极点的 LSZ 约化]
设物理横向二点函数有单粒子简单极点，$Z_{\rm pole}^{(\gamma)}>0$，原始外腿截肢核在该极点光滑，其余部分不含同阶奇异性。
假定 Ward 恒等式及横向偏振条件成立，先在非零 $k^2$ 计算规范固定逆核，再对物理态取留数。
纵向规范等价类不构成额外渐近粒子。
光子外态满足同样的极点逻辑，但只在两个横向物理偏振上取留数。定义物理光子极点权重
\begin{equation}
 \langle\Omega|A_\mu(0)|k,\lambda\rangle
 =\sqrt{Z_{\rm pole}^{(\gamma)}}\,\varepsilon_\mu^{(\lambda)}(k),
 \qquad
 k^2=0,
 \qquad
 k\cdot\varepsilon^{(\lambda)}=0.
 \label{eq:photon-field-pole-overlap-r68}
\end{equation}
若一条光子外腿的连通 Green 函数写成 $\widetilde G_\nu(k;\{q\})=\widetilde D_{\nu\rho}(k)\widetilde\Gamma^\rho(k;\{q\})+\text{非极点项}$，则相应的留数操作为
\begin{equation}
 \boxed{
 \lim_{k^2\to0}
 \varepsilon_\mu^{(\lambda)*}(k)[Z_{\rm pole}^{(\gamma)}]^{-1/2}
 \frac{\mathscr K_{\xi_{\rm gf}}^{\mu\nu}(k)}{\ii\hbar}
 \widetilde G_\nu(k;\{q\})
 =\sqrt{Z_{\rm pole}^{(\gamma)}}\,\varepsilon_\rho^{(\lambda)*}(k)\widetilde\Gamma^\rho(k;\{q\}).}
 \label{eq:photon-lsz-amputation-r68}
\end{equation}
\end{proposition}
\begin{proof}
物理单光子态 $|k,\lambda\rangle$ 只含两个横向偏振，$\lambda=1,2$。其场--粒子重叠可写成

\begin{equation*}
 \langle\Omega|A_\mu(0)|k,\lambda\rangle
 =\sqrt{Z_{\rm pole}^{(\gamma)}}\,\varepsilon_\mu^{(\lambda)}(k),
 \qquad
 k^2=0,
 \qquad
 k\cdot\varepsilon^{(\lambda)}=0.
\end{equation*}
因此精确光子二点函数在横向单光子极点附近具有
\begin{equation*}
 \widetilde D_{\mu\nu}(k)
 =Z_{\rm pole}^{(\gamma)}\,\widetilde D^{(0)}_{\mu\nu}(k)
 +\widetilde D^{\rm reg}_{\mu\nu}(k),
\end{equation*}
其中 $\widetilde D^{\rm reg}_{\mu\nu}$ 在该极点邻域解析或属于连续谱部分。物理四动量表示的自由逆 Maxwell 核由\eqnrefs{eq:photon-physical-momentum-inverse-kernel-dp42} 定义。结合\eqnrefs{eq:photon-prop-kappa-to-q-r13}，可直接相乘得到
\begin{align*}
 \mathscr K_{\xi_{\rm gf}}^{\mu\rho}(k)
 \widetilde D^{(0)}_{\rho\nu}(k)
 &=\left[-\frac{k^2}{\mu_0\hbar^2}
 \left(P_T^{\mu\rho}+\frac1{\xi_{\rm gf}}P_L^{\mu\rho}\right)\right]\\
 &\quad\times
 \left[-\frac{\ii\mu_0\hbar^3}{k^2+\ii0^+}
 \left(P_{T,\rho\nu}+\xi_{\rm gf}P_{L,\rho\nu}\right)\right],\\
 &\xrightarrow[k^2\to0]{}\ii\hbar\,
 (P_T+P_L)^\mu{}_{\nu},\\
 &=\ii\hbar\,\delta^\mu{}_{\nu}.
\end{align*}
第二行使用 $P_T^2=P_T$、$P_L^2=P_L$ 与 $P_TP_L=0$。于是精确传播子的极点部分满足
\begin{equation*}
 [Z_{\rm pole}^{(\gamma)}]^{-1/2}\frac{\mathscr K_{\xi_{\rm gf}}^{\mu\rho}(k)}{\ii\hbar}
 \widetilde D_{\rho\nu}(k)
 \xrightarrow[k^2\to0]{}
 [Z_{\rm pole}^{(\gamma)}]^{1/2}\delta^\mu{}_{\nu}.
\end{equation*}
再与 $\varepsilon_\mu^{(\lambda)*}$ 或 $\varepsilon_\mu^{(\lambda)}$ 收缩，便得到物理光子外态。因而 $(\ii\hbar)^{-1}$ 不是可省略的约定因子，而是由二点函数的归一化由前述约定固定。

不同协变规范只改变含 $k_\mu$ 或 $k_\nu$ 的纵向部分。物理偏振满足
\begin{equation*}
 k_\mu\varepsilon^{\mu(\lambda)}=0,
\end{equation*}
而与守恒物质流连接时 Ward 恒等式给出
\begin{equation*}
 k_\mu\mathcal M^\mu=0.
\end{equation*}
因此 $\varepsilon^\mu\mapsto\varepsilon^\mu+\alpha k^\mu$ 不改变散射矩阵元；光子外腿表示规范等价类中的横向物理态，而不是四势的四个分量各自对应一个粒子。
故对于 $\widetilde G_\nu=\widetilde D_{\nu\rho}\widetilde\Gamma^\rho+\text{非极点项}$，
物理偏振收缩的留数操作留下 $\sqrt{Z_{\rm pole}^{(\gamma)}}\,\varepsilon_\rho^*\widetilde\Gamma^\rho$。
此平方根需计入 LSZ 归一化核，不能在一般方案中从原始截肢核的定义中省去。

\end{proof}


\begin{proposition}[多外腿的逐极点 LSZ 约化]
设各外部不变量具有以上单粒子极点，抽去总四动量守恒分布后，最高阶多重极点可因子化为完整外腿二点函数乘原始核 $\widetilde\Gamma_{\rm raw}$。
其余项至少缺少一个外腿极点，原始核在相应邻域解析。
令 $\widetilde\Gamma_{\rm amp}$ 为原始核乘所有外腿的 $\sqrt{Z_{\rm pole}}$，逐腿留数及壳上投影给出以下振幅。
规范固定核的纵向部分可以出现在中间传播子中；与物理偏振或守恒流收缩后，它不进入散射矩阵元。对每一条外腿重复同一留数操作，连通 Green 函数的最高阶多重极点便分解为外腿极点与一个共同的截肢核，最终得到
\begin{equation}
 \boxed{
 \mathcal M_{fi}
 =\left[\prod_{\rm ext\;fermions}u,\bar u,v,\bar v\right]
 \left[\prod_{\rm ext\;photons}\varepsilon^{(\lambda)}\right]
 \widetilde\Gamma_{\rm amp}.}
 \label{eq:multi-leg-lsz-r68}
\end{equation}
\end{proposition}
\begin{proof}
考虑含 $N_f$ 条费米子外腿和 $N_\gamma$ 条光子外腿的连通精确 Green 函数。若每一个外部不变量都分别接近一个孤立单粒子极点，多变量 Laurent 展开的最高阶极点部分必可分解为

\begin{equation*}
 \widetilde G_{\rm conn}
 =\left(\prod_{a=1}^{N_f}\widetilde S_a\right)
 \left(\prod_{b=1}^{N_\gamma}\widetilde D_b\right)
 \widetilde\Gamma_{\rm raw}
 +\text{至少缺少一个外腿极点的项}.
\end{equation*}
不同外部坐标具有独立 Fourier 共轭动量，因此逐个外部不变量取留数等价于逐条外腿提取局部 Laurent 系数。对每条费米子外腿乘
\begin{equation*}
 [Z_{\rm pole}^{(e)}]^{-1/2}\frac{\slashed p-mc}{\ii\hbar},
\end{equation*}
对每条光子外腿乘
\begin{equation*}
 [Z_{\rm pole}^{(\gamma)}]^{-1/2}\frac{\mathscr K(k)}{\ii\hbar},
\end{equation*}
再投影到相应壳上自旋子和横向偏振后，外腿传播极点被除去，但每条腿留下极点权重的平方根。定义
\[
 \widetilde\Gamma_{\rm amp}
 =\left(\prod_{a=1}^{N_f}\sqrt{Z_{{\rm pole},a}^{(e)}}\right)
  \left(\prod_{b=1}^{N_\gamma}\sqrt{Z_{{\rm pole},b}^{(\gamma)}}\right)
  \widetilde\Gamma_{\rm raw},
\]
并保持 Grassmann 外腿次序以固定费米交换符号，即得到
\begin{equation*}
 \mathcal M_{fi}
 =\left[\prod_{\text{外部费米子}}\text{壳上 }u,\bar u,v,\bar v\right]
 \left[\prod_{\text{外部光子}}\varepsilon^{(\lambda)}\right]
 \widetilde\Gamma_{\rm amp}.
\end{equation*}
所有相互作用顶点的平移不变性共同产生整体
$(2\pi\hbar)^4\delta^{(4)}(P_f-P_i)$；抽去这一总四动量守恒因子后，$\mathcal M_{fi}$ 才进入截面和衰变率公式。
\end{proof}


因此电子与光子的 LSZ 约化共同实现“精确单粒子极点 $\rightarrow$ 渐近物理态 $\rightarrow$ 连通截肢振幅”的映射。
